\documentclass[11pt,a4paper]{article}
\usepackage[utf8]{inputenc}
\usepackage[T1]{fontenc}
\usepackage{fontspec}
\usepackage{xcolor}
\usepackage{graphicx}
\usepackage{geometry}
\usepackage{titlesec}
\usepackage{enumitem}
\usepackage{tcolorbox}
\usepackage{booktabs}
\usepackage{array}
\usepackage{longtable}
\usepackage{hyperref}
\usepackage{fancyhdr}

% 페이지 설정
\geometry{top=25mm, bottom=25mm, left=25mm, right=25mm}

% 한글 폰트 설정
\setmainfont{Noto Sans CJK KR}
\setsansfont{Noto Sans CJK KR}
\setmonofont{JetBrains Mono}[Scale=0.9]

% 색상 정의 (빨간색 제거, 파란색/초록색 사용)
\definecolor{primaryblue}{HTML}{1565C0}
\definecolor{primarygreen}{HTML}{43A047}
\definecolor{accentyellow}{HTML}{FFC107}
\definecolor{textdark}{HTML}{3E2723}
\definecolor{textmedium}{HTML}{5D4037}
\definecolor{bgcream}{HTML}{FFFBF5}
\definecolor{bgblue}{HTML}{E3F2FD}
\definecolor{bggreen}{HTML}{E8F5E9}
\definecolor{bgyellow}{HTML}{FFF8E1}

% 하이퍼링크 색상
\hypersetup{
    colorlinks=true,
    linkcolor=primaryblue,
    urlcolor=primaryblue
}

% 섹션 스타일
\titleformat{\section}
    {\Large\bfseries\color{primaryblue}}
    {\thesection.}{0.5em}{}
\titleformat{\subsection}
    {\large\bfseries\color{primarygreen}}
    {\thesubsection}{0.5em}{}

% 헤더/푸터
\pagestyle{fancy}
\fancyhf{}
\fancyhead[L]{\small\color{textmedium}사과 가격 예측 시뮬레이터 사용설명서}
\fancyhead[R]{\small\color{textmedium}v1.2.1.2}
\fancyfoot[C]{\thepage}
\renewcommand{\headrulewidth}{0.4pt}

% 팁 박스
\newtcolorbox{tipbox}{
    colback=bgyellow,
    colframe=accentyellow,
    boxrule=1pt,
    arc=4pt,
    left=8pt, right=8pt, top=6pt, bottom=6pt,
    title={\textbf{[TIP] 알아두면 좋아요!}},
    fonttitle=\bfseries\color{textdark}
}

% FAQ 박스
\newtcolorbox{faqbox}[1]{
    colback=bgblue,
    colframe=primaryblue,
    boxrule=1pt,
    arc=4pt,
    left=8pt, right=8pt, top=6pt, bottom=6pt,
    title={Q: #1},
    fonttitle=\bfseries\color{primaryblue}
}

% 가이드 박스
\newtcolorbox{guidebox}{
    colback=bggreen,
    colframe=primarygreen,
    boxrule=2pt,
    arc=6pt,
    left=10pt, right=10pt, top=8pt, bottom=8pt
}

\begin{document}

% 제목
\begin{center}
    {\Huge\bfseries\color{primaryblue} 사과 가격 예측 시뮬레이터}\\[8pt]
    {\Large 사용설명서}\\[6pt]
    {\large\color{textmedium} v1.2.1.2}\\[20pt]
\end{center}

% 문서 읽기 순서 안내
\begin{guidebox}
\begin{center}
{\large\bfseries\color{primarygreen} [문서 읽기 순서 안내]}
\end{center}

\begin{center}
\begin{tabular}{ccccccccc}
퀵가이드 & $\rightarrow$ & \colorbox{accentyellow}{\textbf{사용설명서}} & $\rightarrow$ & 이론해설서 & $\rightarrow$ & 활용가이드 & $\rightarrow$ & 수업지도안 \\
(10분) & & (현재, 20분) & & (30분) & & (교사) & & (교사)
\end{tabular}
\end{center}
\end{guidebox}

\vspace{15pt}

% 빠른 시작 가이드
\section*{빠른 시작 가이드 (3단계)}

\begin{center}
\begin{tabular}{|c|l|l|}
\hline
\rowcolor{bggreen}
\textbf{단계} & \textbf{제목} & \textbf{설명} \\
\hline
\cellcolor{primarygreen}\textcolor{white}{\textbf{1}} & 데이터 생성 & 왼쪽 패널에서 \texttt{[데이터 생성]} 버튼 클릭 \\
\hline
\cellcolor{primaryblue}\textcolor{white}{\textbf{2}} & 학습 실행 & 오른쪽 패널에서 \texttt{[학습 실행]} 버튼 클릭 \\
\hline
\cellcolor{accentyellow}\textcolor{textdark}{\textbf{3}} & 추론 모드 & 상단 \texttt{[추론 모드]} 탭 전환 후 슬라이더 조작 \\
\hline
\end{tabular}
\end{center}

\vspace{10pt}

% 이 가이드 사용법
\section*{이 가이드 사용법}

\begin{center}
\begin{tabular}{|c|p{10cm}|}
\hline
\rowcolor{bgblue}
\textbf{단계} & \textbf{내용} \\
\hline
\textbf{1단계} & \textbf{개요 읽기}: 시뮬레이터가 무엇인지, 어떤 개념을 다루는지 파악합니다. \\
\hline
\textbf{2단계} & \textbf{화면 구성 파악}: 어떤 버튼이 어디에 있는지 익힙니다. \\
\hline
\textbf{3단계} & \textbf{따라하기}: 데이터 생성 → 학습 → 예측 순서로 직접 실습합니다. \\
\hline
\textbf{4단계} & \textbf{심화 학습}: 파라미터를 바꿔가며 결과 변화를 관찰합니다. \\
\hline
\end{tabular}
\end{center}

\tableofcontents
\newpage

%--------------------------------------------------
\section{개요}
%--------------------------------------------------

이 시뮬레이터는 \textbf{선형 회귀(Linear Regression)}의 원리를 학습하기 위한 교육용 도구입니다. 
사과의 \textbf{당도}와 \textbf{무게}를 \textbf{품질 점수}로 통합하여 \textbf{가격}을 예측하는 과정을 시각적으로 체험할 수 있습니다.

\begin{tipbox}
\textbf{왜 품질 점수로 통합하나요?}

우리가 배우는 ``단순 선형 회귀''는 입력 변수가 \textbf{1개}여야 합니다. 
당도와 무게, 두 가지 정보를 하나로 합친 ``품질''이라는 대표 점수를 만들어서 
$Y = wX + b$ 형태의 간단한 직선 관계를 학습합니다.
\end{tipbox}

\subsection{시스템 요구사항}

\textbf{핵심 요약}: 크롬 브라우저만 있으면 바로 실행 가능합니다!

\begin{center}
\begin{tabular}{ll}
\toprule
\textbf{항목} & \textbf{요구사항} \\
\midrule
웹 브라우저 & Chrome, Edge, Firefox (권장: Chrome) \\
화면 해상도 & 1280×720 이상 \\
인터넷 연결 & 초기 로딩 시 필요 (폰트, 라이브러리) \\
\bottomrule
\end{tabular}
\end{center}

\subsection{핵심 특징 (v1.2.4.1)}

\begin{itemize}[leftmargin=*]
    \item \textbf{품질 점수 통합}: 당도 + 무게 → 0\textasciitilde100점
    \item \textbf{단순 선형 회귀}: 가격 = $w \times$ 품질 $+ b$
    \item \textbf{2D 시각화}: X축(품질), Y축(가격) 산점도
    \item \textbf{오차선 분석}: 예측과 실제의 차이를 시각화
    \item \textbf{학습 과정 재생}: 애니메이션으로 학습 과정 관찰
\end{itemize}

%--------------------------------------------------
\section{화면 구성}
%--------------------------------------------------

화면은 크게 \textbf{3개 영역}으로 나뉩니다.

\begin{center}
\begin{tabular}{|c|c|c|}
\hline
\rowcolor{bgblue}
\textbf{왼쪽 패널 (240px)} & \textbf{중앙 차트 영역} & \textbf{오른쪽 패널 (260px)} \\
\hline
데이터 생성 & 손실 그래프 & 학습 설정 \\
품질 공식 설명 & 산점도 + 회귀선 & (학습률, 에포크) \\
\hline
\end{tabular}
\end{center}

\subsection{모드 전환}

상단 탭으로 두 가지 모드를 전환합니다.

\begin{center}
\begin{tabular}{|l|p{8cm}|}
\hline
\rowcolor{bggreen}
\textbf{모드} & \textbf{설명} \\
\hline
\textbf{학습 모드} & 데이터를 생성하고, AI가 패턴을 학습합니다. \\
\hline
\textbf{추론 모드} & 학습된 모델로 새로운 사과의 가격을 예측합니다. \\
\hline
\end{tabular}
\end{center}

\begin{tipbox}
\textbf{학습과 추론의 관계}

``공부(학습)를 해야 시험(예측)을 볼 수 있다''고 생각하세요! 
학습 모드에서 만들어진 ``회귀식''이 추론 모드의 예측 공식으로 그대로 전달됩니다.
\end{tipbox}

%--------------------------------------------------
\section{품질 점수 계산}
%--------------------------------------------------

당도와 무게를 하나의 \textbf{품질 점수(0\textasciitilde100)}로 통합합니다.

\[
\text{품질} = \text{당도점수}(50\%) + \text{무게점수}(50\%)
\]

\begin{center}
\begin{tabular}{lll}
\toprule
\textbf{항목} & \textbf{계산식} & \textbf{범위} \\
\midrule
당도 점수 & $(\text{당도} - 10) / 8 \times 50$ & 10\textasciitilde18 Brix → 0\textasciitilde50점 \\
무게 점수 & $(\text{무게} - 150) / 200 \times 50$ & 150\textasciitilde350g → 0\textasciitilde50점 \\
\midrule
총 품질 & 당도점수 + 무게점수 & 0\textasciitilde100점 \\
\bottomrule
\end{tabular}
\end{center}

\textbf{예시}:
\begin{itemize}
    \item 당도 14 Brix: $(14 - 10) / 8 \times 50 = 25$점
    \item 무게 250g: $(250 - 150) / 200 \times 50 = 25$점
    \item 총 품질: $25 + 25 = 50$점
\end{itemize}

%--------------------------------------------------
\section{기능별 사용법}
%--------------------------------------------------

\subsection{데이터 생성}

\begin{center}
\begin{tabular}{lll}
\toprule
\textbf{파라미터} & \textbf{범위} & \textbf{권장값} \\
\midrule
데이터 개수 & 5\textasciitilde100개 & 20\textasciitilde50개 \\
노이즈 수준 & 5\% / 10\% / 20\% & 10\% \\
\bottomrule
\end{tabular}
\end{center}

\begin{tipbox}
\textbf{노이즈란?}

현실 세계의 데이터는 완벽하지 않습니다. 같은 품질의 사과라도 가격이 조금씩 다를 수 있죠.
노이즈는 이런 ``현실의 불확실성''을 표현합니다. 노이즈가 높을수록 예측이 어려워집니다.
\end{tipbox}

\subsection{모델 학습}

\begin{center}
\begin{tabular}{llp{6cm}}
\toprule
\textbf{파라미터} & \textbf{권장값} & \textbf{설명} \\
\midrule
학습률 & 0.01 & 한 번에 얼마나 크게 움직일지 결정. 너무 크면 불안정, 너무 작으면 느림 \\
에포크 & 500 & 전체 데이터를 몇 번 반복 학습할지. 많을수록 정확하지만 시간 소요 \\
\bottomrule
\end{tabular}
\end{center}

\subsection{학습 결과 분석 (3열 패널)}

학습 완료 후 차트 아래에 3개 패널이 표시됩니다.

\begin{center}
\begin{tabular}{|l|p{8cm}|}
\hline
\rowcolor{bgblue}
\textbf{패널} & \textbf{내용} \\
\hline
학습 결과 & 에포크, MSE 손실, R² 점수, 회귀식 \\
\hline
오차 분석 & 오차선 토글, 평균/최대 오차, MSE \\
\hline
학습 재생 & 재생 속도 조절, 재생/정지 버튼 \\
\hline
\end{tabular}
\end{center}

\subsection{예측 (추론 모드)}

슬라이더를 움직이면 \textbf{자동으로} 예측 결과가 표시됩니다.

\begin{enumerate}
    \item 상단에서 \texttt{[추론 모드]} 탭 클릭
    \item 왼쪽 패널에서 당도, 무게 슬라이더 조절
    \item 오른쪽 패널에 예측 가격이 즉시 표시됨
    \item 차트에 수직선/수평선/예측점이 시각화됨
\end{enumerate}

%--------------------------------------------------
\section{시각화 요소}
%--------------------------------------------------

\begin{center}
\begin{tabular}{ll}
\toprule
\textbf{요소} & \textbf{설명} \\
\midrule
X축 & 품질 점수 (0\textasciitilde100) \\
Y축 & 가격 (원) \\
데이터 점 & 1.5px 크기, 녹색(저품질)→붉은색(고품질) 색상 변화 \\
회귀선 & 파란색 직선 (학습된 예측 모델) \\
오차선 & 점선, 오차 크기에 따라 색상 변화 \\
예측점 & 노란 원 + 수직/수평 점선 \\
\bottomrule
\end{tabular}
\end{center}

%--------------------------------------------------
\section{학습 결과 해석}
%--------------------------------------------------

\subsection{회귀식 이해}

학습 후 표시되는 회귀식 예시:
\[
\text{가격} = 30.5 \times \text{품질} + 985
\]

\begin{itemize}
    \item \textbf{가중치 30.5}: 품질이 1점 오르면 가격이 약 30원 상승
    \item \textbf{편향 985}: 품질이 0점일 때의 기본 가격
    \item 품질 100점이면: $30.5 \times 100 + 985 = 4,035$원 예측
\end{itemize}

\subsection{R² (결정계수) 해석}

\begin{tipbox}
\textbf{R²가 뭔가요?}

``모델이 데이터를 얼마나 잘 설명하는가''를 나타내는 점수입니다.
100점 만점이라고 생각하세요. 80점 이상이면 ``꽤 믿을만하다!''라는 뜻입니다.
\end{tipbox}

\begin{center}
\begin{tabular}{lll}
\toprule
\textbf{R² 범위} & \textbf{상태} & \textbf{의미} \\
\midrule
$> 0.8$ & [양호] & 모델이 데이터를 잘 설명함 \\
$0.5 \sim 0.8$ & [보통] & 어느 정도 설명하지만 개선 여지 있음 \\
$\leq 0.5$ & [부족] & 모델 개선 필요 \\
\bottomrule
\end{tabular}
\end{center}

\subsection{MSE (평균 제곱 오차)}

\begin{tipbox}
\textbf{MSE가 뭔가요?}

``예측이 틀린 정도''를 숫자로 나타낸 것입니다. 
작을수록 예측이 정답에 가깝습니다. 0에 가까울수록 좋아요!
\end{tipbox}

%--------------------------------------------------
\section{문제 해결 (FAQ)}
%--------------------------------------------------

\begin{faqbox}{손실 그래프가 위아래로 요동쳐요}
\textbf{A}: 학습률이 너무 큽니다. 
\texttt{학습률}을 \textbf{0.01 이하}로 낮춰보세요. 
학습률이 크면 ``한 번에 너무 많이 움직여서'' 정답을 지나쳐버립니다.
\end{faqbox}

\begin{faqbox}{R² 점수가 음수예요}
\textbf{A}: 모델이 ``평균으로 찍는 것보다 못한'' 예측을 하고 있습니다.
\begin{itemize}
    \item 데이터 개수를 늘려보세요 (30개 이상)
    \item 에포크를 늘려보세요 (500회 이상)
    \item 노이즈를 낮춰보세요 (10\% 이하)
\end{itemize}
\end{faqbox}

\begin{faqbox}{추론 모드에서 예측이 안 돼요}
\textbf{A}: 학습을 먼저 완료해야 합니다!
학습 모드에서 \texttt{[학습 실행]} 버튼을 클릭하고, 
``학습 완료'' 메시지가 나온 후에 추론 모드로 전환하세요.
\end{faqbox}

\begin{faqbox}{학습이 너무 오래 걸려요}
\textbf{A}: 에포크 수를 줄여보세요. 
500회 정도면 대부분의 경우 충분합니다.
\end{faqbox}

%--------------------------------------------------
\section{용어 해설 (Glossary)}
%--------------------------------------------------

\begin{center}
\begin{longtable}{lp{9cm}}
\toprule
\textbf{용어} & \textbf{설명} \\
\midrule
\endhead

선형 회귀 & 직선 관계를 찾는 방법. $Y = wX + b$ 형태의 수식을 학습 \\
\midrule
가중치 (w) & 기울기. X가 1 증가할 때 Y가 얼마나 변하는지 \\
\midrule
편향 (b) & Y절편. X가 0일 때의 Y값 \\
\midrule
학습률 & 한 번에 가중치를 얼마나 조정할지. 보폭 크기와 비슷 \\
\midrule
에포크 & 전체 데이터를 한 번 학습하는 것. 100 에포크 = 100번 반복 \\
\midrule
MSE & Mean Squared Error. 오차의 제곱 평균. 작을수록 좋음 \\
\midrule
R² & 결정계수. 모델의 설명력. 1에 가까울수록 좋음 \\
\midrule
손실 (Loss) & 예측과 실제의 차이. 학습의 목표는 손실 최소화 \\
\midrule
경사 하강법 & 손실을 줄이는 방향으로 가중치를 조금씩 조정하는 방법 \\
\midrule
오차선 (잔차) & 실제값과 예측값의 차이를 시각화한 선 \\
\midrule
노이즈 & 데이터의 불확실성. 현실 세계의 랜덤한 변동 \\
\midrule
정규화 & 서로 다른 범위의 값을 동일한 기준으로 맞추는 것 \\
\midrule
품질 점수 & 당도+무게를 0\textasciitilde100점으로 통합한 점수 \\
\midrule
회귀선 & 학습된 직선. 데이터의 패턴을 대표하는 선 \\
\bottomrule
\end{longtable}
\end{center}

%--------------------------------------------------
\section*{참고자료}
%--------------------------------------------------

\begin{center}
\begin{tabular}{clll}
\toprule
\textbf{\#} & \textbf{자료명} & \textbf{유형} & \textbf{비고} \\
\midrule
1 & 사과가격예측\_선형회귀\_시뮬레이터\_v1.2.4.1.html & 시뮬레이터 & 최신 버전 \\
2 & 사과가격예측\_선형회귀\_시뮬레이터\_개발일지\_v1.2.4.1.md & 개발일지 & 기술 상세 \\
3 & 문서버전관리가이드\_v1.2.2 & 지침 & 버전 관리 \\
\bottomrule
\end{tabular}
\end{center}

%--------------------------------------------------
\section*{생성/수정 이력}
%--------------------------------------------------

\begin{center}
\begin{tabular}{llllll}
\toprule
\textbf{버전} & \textbf{날짜} & \textbf{시간} & \textbf{변경 수준} & \textbf{변경 내용} & \textbf{작성자} \\
\midrule
v1.0.0.0 & 2026-01-31 & 14:00 & 최초 & 초안 작성 & Claude \\
v1.1.0.0 & 2026-01-31 & 15:40 & 마이너 & v1.2.4.1 기준 업데이트 & Claude \\
v1.2.0.0 & 2026-02-01 & 04:30 & 마이너 & 교육용 그림 삽입 & Claude \\
v1.2.0.4 & 2026-02-01 & 15:00 & 패치 & 읽기순서, 빠른시작 추가 & Claude \\
v1.2.1.0 & 2026-02-01 & 18:00 & 마이너 & 초보자 개선, FAQ, 용어집 & Claude \\
\textbf{v1.2.1.2} & \textbf{2026-02-01} & \textbf{18:40} & \textbf{패치} & \textbf{LaTeX 변환, 이모지→텍스트, 색상 변경} & \textbf{Claude} \\
\bottomrule
\end{tabular}
\end{center}

\vfill
\begin{center}
\textit{사과 가격 예측 선형 회귀 시뮬레이터 사용설명서 v1.2.1.2}
\end{center}

\end{document}
